% Part: computability
% Chapter: recursive-functions
% Section: pr-functions

\documentclass[../../../include/open-logic-section]{subfiles}

\begin{document}

\olfileid{cmp}{rec}{prf}
\olsection{原始递归函数}

让我们再次回顾如何利用原始递归和复合从已有函数定义新函数。

\begin{defn}
  \ollabel{defn:primitive-recursion} 设 $f$ 是一个 $k$ 元函数（$k\ge 1$），$g$ 是一个 $(k+2)$ 元函数。\emph{由 $f$ 和 $g$ 经原始递归定义的}函数是 $(k+1)$ 元函数~$h$，定义为
  \begin{align*}
    h(x_0,\dots,x_{k-1},0) & =  f(x_0,\dots,x_{k-1}) \\
    h(x_0,\dots,x_{k-1},y+1) & = g(x_0,\dots,x_{k-1}, y, h(x_0,\dots,x_{k-1}, y))
  \end{align*}
\end{defn}

\begin{defn}
  \ollabel{defn:composition} 设 $f$ 是一个 $k$ 元函数，且 $g_0$、\dots、$g_{k-1}$ 是 $k$ 个 $n$ 元函数。\emph{由 $f$ 与 $g_0$、\dots、~$g_{k-1}$ 经复合定义的}函数是 $n$ 元函数~$h$，定义为
  \[
  h(x_0, \dots, x_{n-1}) =
  f(g_0(x_0, \dots, x_{n-1}), \dots, g_{k-1}(x_0, \dots, x_{n-1})).
  \]
\end{defn}

除了 $\Succ$ 以及对于每个自然数 $n$ 和 $i < n$ 的投影函数
\[
\Proj{n}{i}(x_0,\dots,x_{n-1}) = x_i,
\]
之外，我们还将函数 $\Zero(x) = 0$ 纳入原始递归函数。

\begin{defn}
  原始递归函数集是从 $\Nat^n$ 到 $\Nat$ 的函数集，由以下条款归纳地定义：
  \begin{enumerate}
  \item $\Zero$ 是原始递归的。
  \item $\Succ$ 是原始递归的。
  \item 每个投影函数 $\Proj{n}{i}$ 都是原始递归的。
  \item 如果 $f$ 是一个 $k$ 元原始递归函数，且 $g_0$,
    \dots,~$g_{k-1}$ 是 $n$ 元原始递归函数，那么 $f$ 与 $g_0$, \dots,~$g_{k-1}$ 的复合是原始递归的。
  \item 如果 $f$ 是一个 $k$ 元原始递归函数，且 $g$ 是一个
    $k+2$ 元原始递归函数，那么由 $f$ 和 $g$ 经原始递归定义的函数是原始递归的。
\end{enumerate}
\end{defn}

\begin{explain}
更简洁地说，原始递归函数集是包含 $\Zero$、$\Succ$ 以及投影函数~$\Proj{n}{j}$，并且在复合与原始递归下封闭的最小集合。

描述原始递归函数集的另一种方式是按“阶段”来定义它。令 $S_0$ 表示初始函数集：$\Zero$、$\Succ$ 和投影函数。这些是第~$0$ 阶段的原始递归函数。一旦定义了阶段 $S_i$，令 $S_{i+1}$ 为对 $S_i$ 中的函数应用一次复合或原始递归所得到的所有函数的集合。那么
\[
S = \bigcup_{i \in \Nat} S_i
\]
就是所有原始递归函数的集合。
\end{explain}

我们来验证 $\Add$ 是原始递归函数。

\begin{prop}
  加法函数 $\Add(x,y) = x+y$ 是原始递归的。
\end{prop}

\begin{proof}
我们已经有了一个关于 $\Add$ 的原始递归定义，它通过两个函数 $f$ 和~$g$ 来定义，并且符合
\olref{defn:primitive-recursion} 的格式：
\begin{align*}
  \Add(x_0, 0) & = f(x_0) = x_0\\
  \Add(x_0, y+1) & = g(x_0, y, \Add(x_0, y)) =
  \Succ(\Add(x_0, y))
\end{align*}
所以，只要 $f$ 和 $g$ 是原始递归的，$\Add$ 就是原始递归的。$f(x_0) = x_0 = \Proj{1}{0}(x_0)$，而投影函数
视为原始递归函数，因此 $f$ 是原始递归的。函数 $g$ 是三元函数 $g(x_0, y, z)$，定义为
\[
g(x_0, y, z) = \Succ(z).
\]
这还不能说明 $g$ 是原始递归的，因为 $g$ 和 $\Succ$ 并不完全相同：$\Succ$ 是一元的，而 $g$ 必须是三元的。但我们可以“形式上”通过复合将 $g$ 定义为
\[
g(x_0, y, z) = \Succ(\Proj{3}{2}(x_0, y, z))
\]
由于 $\Succ$ 和 $\Proj{3}{2}$ 视为原始递归函数，而 $g$ 可以由原始递归函数复合定义，因此 $g$ 也是原始递归的。
\end{proof}

\begin{prop}
  \ollabel{prop:mult-pr}
  乘法函数 $\Mult(x,y) = x \cdot y$ 是原始递归的。
\end{prop}

\begin{proof}
  习题。
\end{proof}

\begin{prob}
  通过说明 $\Mult$ 的原始递归定义可以写成 \olref[cmp][rec][prf]{defn:primitive-recursion} 所要求的形式，并证明相应的函数 $f$ 和~$g$ 是原始递归的，来证明 \olref[cmp][rec][prf]{prop:mult-pr}。
\end{prob}

\begin{ex}
这是我们关于原始递归定义的第一个例子：
\begin{align*}
h(0) & =  1 \\
h(y+1) & =  2 \cdot h(y).
\end{align*}
这个函数无法直接套用 \olref{defn:primitive-recursion} 所要求的形式，因为 $k=0$。此外，定义中还涉及常数 $1$ 和 $2$。为了解决第一个问题，我们引入一个哑自变元并定义函数~$h'$：
\begin{align*}
h'(x_0, 0) & =  f(x_0) = 1 \\
h'(x_0, y+1) & =  g(x_0, y, h'(x_0, y)) = 2 \cdot h'(x_0, y).
\end{align*}
函数 $f(x_0) = 1$ 可以通过 $\Succ$ 和 $\Zero$ 的复合来定义：$f(x_0) = \Succ(\Zero(x_0))$。函数 $g$ 可以通过 $g'(z) = 2 \cdot z$ 与投影函数的复合来定义：
\begin{align*}
g(x_0, y, z) & = g'(\Proj{3}{2}(x_0, y, z))
\intertext{而 $g'$ 又可以进而通过复合定义为：}
g'(z) & = \Mult(g''(z), \Proj{1}{0}(z))
\intertext{以及}
g''(z) & = \Succ(f(z)),
\end{align*}
这里的 $f$ 就是上面定义的：$f(z) = \Succ(\Zero(z))$。现在我们有了~$h'$，可以再次使用复合来定义 $h(y) =
h'(\Proj{1}{0}(y),\Proj{1}{0}(y))$。这表明 $h$ 可以通过一系列复合和原始递归从基本函数定义出来，因此 $h$~是原始递归的。
\end{ex}

\end{document}